\documentclass[11pt]{article}
\usepackage[T1]{fontenc}
\usepackage{makeidx}
\usepackage{hyperref}
\hypersetup{colorlinks=true,linkcolor=blue,citecolor=blue,urlcolor=blue}
\makeindex

\title{Spacing, Glue, and Friends}
\author{S. Liang}
\date{}

\begin{document}
\maketitle

\section{Why spacing matters}
\label{sec:why}

Good spacing\index{spacing|(} is invisible: the reader notices nothing, and
that is the point. Spacing rests on three ideas. First, the
\index{space!interword}interword space is glue, not a fixed blank, so each
line can breathe a little; the details of vertical rhythm wait until
Section~\ref{sec:rhythm} on page~\pageref{sec:rhythm}. Second,
sentence endings\index{space!sentence}
traditionally get extra room in English typography. Third, none of this works
without the paragraph builder distributing the slack evenly.

The synonymous term \index{leading|see{spacing}}leading belongs to the days
of metal type, and readers looking for it should consult spacing instead.
Likewise \index{letterspacing|see{spacing}}letterspacing is covered here.

\section{Glue in detail}

Glue has three parts: a natural size, a stretch, and a shrink. When \TeX{}
sets a paragraph, it chooses line breaks and then sets each line by
stretching or shrinking the glue on it. A line stretched too far is
underfull; a line shrunk too far sticks into the margin. The badness number
quantifies the complaint, and the demerits total decides between candidate
breakpoints.

Consider a narrow column of justified text. Every line must reach both
margins, so the glue does real work on each one. Hyphenation helps by
splitting long words, which gives the breaker more candidate endpoints and
keeps the glue closer to its natural size. Without hyphenation, narrow
columns produce rivers of white space and frequent overfull boxes.

\section{Vertical rhythm}
\label{sec:rhythm}

Pages need rhythm as well as lines. The baseline skip separates consecutive
lines, and the paragraph skip separates thoughts. Headings interrupt the
rhythm deliberately: extra space above a heading groups what follows under
it. Footnotes borrow space from the bottom of the page, and floats negotiate
for whole regions at the top or bottom.

Keeping this rhythm steady across pages is mostly a matter of consistent
glue: if every heading, list, and display uses the same set of skips, the
page grid holds together.

One more device deserves mention: the strut, an invisible rule that forces
a minimum line height wherever alignment matters, as in tables and display
mathematics. Struts do not stretch or shrink; they simply guarantee room.
With struts in place, even a page of mixed headings, lists, and displays
keeps its even color. That concludes our tour of
spacing\index{spacing|)} as a full-page concern.

\end{document}
